% CAPITOLO 5 — EXCEPTIONS
\chapter{Exceptions}
\label{cap:exceptions}

Il progetto introduce un piccolo insieme di eccezioni applicative per separare
gli errori di \emph{validazione del dominio} (regole di business) dalla
gestione tecnica dell'infrastruttura. La gerarchia è organizzata attorno alla
classe base \texttt{BusinessValidationException}; ogni eccezione è \emph{checked}
e viene catturata a livello di \texttt{Boundary} per la presentazione del
messaggio all'utente, senza mai propagare stack infratilistici fuori dai
confini BCE.

\section{BusinessValidationException}
\texttt{BusinessValidationException} estende \texttt{java.lang.Exception}
ed è la radice delle violazioni di regola di business nel progetto. Viene
sollevata ogni volta che l'input dell'utente (o lo stato del dominio) viola
una regola applicativa, ad esempio:
\begin{itemize}
    \item \texttt{OrdinaProdottoController} la solleva quando l'autorizzazione
    del pagamento è rifiutata (estensione \textbf{6a}) o le scorte sono
    insufficienti (estensione \textbf{3a});
    \item la validazione dei campi carta (\texttt{validaCampi}) la usa per
    segnalare campi mancanti o CVV non valido;
    \item il buono promozionale non valido (inesistente, scaduto, già usato,
    venditore errato) ritorna \texttt{BusinessValidationException}.
\end{itemize}
Gradualmente le Boundary la catturano e mostrano il messaggio in una
\texttt{Label} inline (nessuna view fittizia), mantenendo l'utente nella
schermata corrente.

\section{AutenticazioneException}
Sottoclasse di \texttt{BusinessValidationException} che racchiude gli errori
di \textbf{autenticazione} e \textbf{registrazione}: credenziali errate,
utente o password non valide, violazioni sulla creazione del profilo. È
usata dal controllare \texttt{AutenticazioneController} per uniformare i
messaggi di login.

\section{InvalidStateTransitionException}
Sottoclasse di \texttt{BusinessValidationException} usata dalla macchina a
stati dell'\texttt{Ordine} (BR-04): quando si tenta una transizione di stato
non prevista (es. passare direttamente da \texttt{CREATED} a
\texttt{DELIVERED}), \texttt{Ordine.cambiaStato()} solleva
\texttt{InvalidStateTransitionException}, bloccando l'operazione illegale e
conservando l'integrità del modello.

\section{Gestione nelle Boundary}
Tutte le eccezioni applicative sono catturate nei \texttt{handler} delle
view JavaFX (es. \texttt{catcher} \texttt{BusinessValidationException} nel
\texttt{setOnAction} della \texttt{PaymentView}): il messaggio viene mostrato
inline nella \emph{Label} di stato ed è fornito dal costruttore
dell'eccezione; ciò garantisce coerenza con la modellazione BCE (le Entity
non dipendono da GUI) e con lo State Diagram a \emph{schermate}.

\textbf{Nota:} gli errori di infrastruttura invece (es. \texttt{SQLException})
sono gestiti dai DAO e non si propagano alle view, come documentato nel
Capitolo~\ref{cap:persistenza}; il Quality Gate impone test automatizzati
per tutti i rami di errore (Capitolo~\ref{cap:testing}).